\scanpage{67}
Dans le cas A, on appellera \uline{repère} \uline{du} graphe \inserted{$P_M=$} $(\unsure{S|B}_M,A_M)$ un élément de $\overrightarrow A_M$, \unsure{ie.|i.e.}
\[
 \operatorname{Rep}(P_M)\overset{\mathrm{def}}{=}\overrightarrow A_M\quad(\text{cas }2.1_{\Lambda}=0),\tag*{(176)}
\]
dans le cas B), on définit $\operatorname{Rep}(P_M)$ \deleted{\unsure{(174)}} par (174). Dans tous les cas, l'ens des repères \unsure{de.|de|abs.} $P_M$ est un torseur sous $G_M$.

\uline{Relations entre les données} $\deleted{S}A_M,F_M$. Je voudrais m'assurer que ces données sur $S_M$
\[
 \left[\begin{gathered}A_M\subset\mathfrak P_2(S_M)\\A_M\longrightarrow\unsure{\pi|\Pi}\end{gathered}\right],\
 \left[\begin{gathered}F_M\subset\mathfrak P_3(S_M)\\F_M\longrightarrow\pi\end{gathered}\right]\tag*{(177)}
\]
% slanted
\margin{$A_M\mathbin{-}\pi$, $F_M\mathbin{-}\pi$ sont \unsure{interprétés|interprétées} ici comme des partitions de $A_M,F_M$ --- \unsure{le|la} \deleted{\unsure{but}} de \unsure{l'ens d'indices|l'un d'eux} $\pi$ étant pour le moment, \unsure{irrelevant}\ldots}
se déterminent mutuellement, En tous cas, dans le cas A\deleted{\unsure{B)} $(2.1_\Lambda\ne0)$}\inserted{comme B}, on a
\[
 A_M=\{a\in\mathfrak P_2(S_M)\mid\exists f\in F_M\deleted{\unsure{\}}},f\supset a\}\tag*{(177)}
\]
(168, \inserted{170,} 171), et d'autre part je \deleted{\unsure{dis}} \inserted{voudrais} \deleted{que}\inserted{\uuline{ds le cas} B $(2.1_\Lambda\ne0)$}
\[
 F_M=\{f\in\mathfrak P_3(S_M)\mid\forall a\in\mathfrak P_2(f),\text{ on a }a\in A_M\}\tag*{\deleted{(178)}}
\]
% slanted
\margin{relation fausse !!}
On a bien $F_M\subset\{\ \}$, par (168), inversement \deleted{\unsure{si on|a-t-on}} a-t-on $F_M\supset\{\ \}$ ?

\uuline{C'est ultra faux} même sur $\Lambda=\mathbb Z/n\mathbb Z$ ou $\Lambda=\mathbb Z$, cela signifierait que si $x,y,z\in M$ sont tels que $x\wedge y,y\wedge z,z\wedge x$ sont dans $\det M^*$, alors on peut \deleted{\unsure{que}} trouver
% archive sheet number: 66
